\documentclass[11pt]{article}

\usepackage[margin=1in]{geometry}
\usepackage{parskip}
\usepackage{amsmath, amssymb, amsthm}
\usepackage{booktabs}
\usepackage{hyperref}

\title{Summary statistics from scratch}
\author{Group N}
\date{}

\begin{document}
\maketitle

\begin{abstract}
We present \texttt{stats\_from\_scratch}, a small Python library
that implements the mean, the median, the population and sample
standard deviations, and the three quartiles directly from their
definitions. Every function is checked against the matching function
in the standard library \texttt{statistics} module; we treat the
standard library as a reference implementation in the tests, not in
the library code. The library follows the Diataxis framework for
documentation~\cite{procida2017} and the statistics chapter of the
open-source textbook \emph{Python for
Mathematics}~\cite{knight2024}; we mention the wider statistical
context in~\cite{rice2007}.
\end{abstract}

\section{Introduction}

The summary statistics of a sample are the most basic description of
its location and spread. For a sample
\(x_1, x_2, \dots, x_n\) the arithmetic mean is
\(\bar{x} = (1/n) \sum_i x_i\), the median is the middle of the
sorted sample, and the standard deviation is the square root of the
mean squared deviation from the mean. The two common variants of the
standard deviation differ in the divisor: \(n\) for the population
version and \(n - 1\) for the sample version~\cite{rice2007}.

The standard library \texttt{statistics} module already implements
all of these. We wanted a library that we could read and understand
in full, with the formulas in plain sight, and we treat the standard
library as the reference implementation in our tests.

We present a five-function library,
\texttt{stats\_from\_scratch}; every function is implemented from
its definition; and every function is checked against the matching
function in the standard library \texttt{statistics} module.

\section{The library}

The library lives in the single module
\texttt{stats\_from\_scratch.py} and follows the statistics chapter
of \emph{Python for Mathematics}~\cite{knight2024}.

\texttt{mean(values)} returns
\(\sum_i x_i / n\).

\texttt{median(values)} sorts the sample,
returns the middle value for an odd count, and the average of the
two middle values for an even count.

\texttt{population\_standard\_deviation(values)} and
\texttt{sample\_standard\_deviation(values)} return the square root
of the mean squared deviation with divisors \(n\) and \(n - 1\)
respectively.

\texttt{quartiles(values)} returns the three
cut points \(Q_1\), \(Q_2\), \(Q_3\), using the 'exclusive' method
of the standard library: the \(k\)-th cut is at the fractional
1-indexed position \(k (n + 1) / 4\) of the sorted sample, with
linear interpolation between the surrounding values. We picked this
method so the library agrees with \texttt{statistics.quantiles(data,
n=4)} exactly.

\section{Worked example}

We illustrate the library on the sample \([1, 2, 3, 4, 5]\). The
mean is \(3\), the median is \(3\), and the three quartiles are
\((1.5, 3, 4.5)\), as expected. The population standard deviation is
\(\sqrt{2} \approx 1.414\) and the sample standard deviation is
\(\sqrt{2.5} \approx 1.581\). The values match
\texttt{statistics.mean}, \texttt{statistics.median},
\texttt{statistics.pstdev}, \texttt{statistics.stdev}, and
\texttt{statistics.quantiles(data, n=4)} to within floating-point
precision (Table~\ref{tab:demo}).

\begin{table}[h]
\centering
\begin{tabular}{l l l}
\toprule
Quantity & Library & \texttt{statistics} \\
\midrule
mean    & \(3.0\)              & \(3\) \\
median  & \(3\)                & \(3\) \\
pstdev  & \(\sqrt{2}\)         & \(\sqrt{2}\) \\
stdev   & \(\sqrt{2.5}\)       & \(\sqrt{2.5}\) \\
quartiles & \((1.5, 3, 4.5)\)  & \([1.5, 3.0, 4.5]\) \\
\bottomrule
\end{tabular}
\caption{\textbf{Agreement with the standard library on
$[1, 2, 3, 4, 5]$.} Every function in the library returns a value
that matches \texttt{statistics} to within floating-point
precision.}
\label{tab:demo}
\end{table}

\section{Discussion}

The standard library already covers
this material in full. We wrote the library because we wanted to
know what the formulas were doing, and the act of writing them out
in clean form was itself instructive.

There are several methods in
use~\cite{rice2007}. We picked the 'exclusive' method so that the
library agrees with the standard library's
\texttt{statistics.quantiles}. A natural extension would be to
expose a ``method'' keyword so that other methods can be selected.

The library handles only one-dimensional
samples. A natural extension is the two-sample case (correlation,
covariance, Pearson and Spearman coefficients), each of which has a
short definition and an obvious unit-test target in
\texttt{statistics} or \texttt{scipy.stats}.

\section{Conclusion}

We have presented \texttt{stats\_from\_scratch}, a small library
that implements the basic summary statistics from their definitions
and agrees with the standard library \texttt{statistics} module on
every test case we tried.

\bibliographystyle{plain}
\bibliography{references}

\end{document}
